% Part: first-order-logic
% Chapter: beyond
% Section: modal-logics

\documentclass[../../../include/open-logic-section]{subfiles}

\begin{document}

\olfileid{fol}{byd}{mod}

\olsection{মোডাল যুক্তিবিদ্যা}

নিচের শর্তাধীন বাক্যটির উদাহরণ বিবেচনা করো:
\begin{quote}
  জেরেমি যদি ওই ঘরে একা থাকে, তাহলে সে মাতাল ও নগ্ন এবং চেয়ারগুলির
  উপর নাচছে।
\end{quote}
এটি এমন একটি শর্তাধীন দাবির উদাহরণ যা বস্তুগতভাবে সত্য হলেও বিভ্রান্তিকর
হতে পারে। কারণ বাক্যটি যেন ইঙ্গিত করে যে পূর্বপক্ষ ও সিদ্ধান্তের মধ্যে
আরও জোরালো কোনো যোগ আছে---শুধু পূর্বপক্ষটি মিথ্যা অথবা অনুপক্ষটি সত্য,
এটুকুই নয়। অর্থাৎ, শব্দচয়নটি ইঙ্গিত করে যে দাবিটি কেবল এই বিশেষ
জগতে সত্য নয় (এখানে জেরেমি ঘরে একা না থাকার কারণে সেটি তুচ্ছভাবেই
সত্য হতে পারে), বরং পূর্বপক্ষটি সত্য \emph{হলে} অনুপক্ষটিও সত্য
\emph{হতো}। অন্যভাবে বললে, দাবিটির অর্থ এমন নেওয়া যায় যে এটি শুধু
এই জগতেই নয়, যেকোনো “সম্ভাব্য” জগতেও সত্য; অথবা, এই বিশেষ জগতে
নিছক সত্য হওয়ার বিপরীতে এটি \emph{আবশ্যিকভাবে} সত্য।

এই ধরনের আবশ্যিকতাকে অর্থপূর্ণ করতেই মোডাল যুক্তিবিদ্যা নির্মিত।
সাধারণ প্রস্তাবনামূলক যুক্তিবিদ্যায় একটি বক্স অপারেটর যোগ করলে মোডাল
প্রস্তাবনামূলক যুক্তিবিদ্যা পাওয়া যায়; অর্থাৎ, $!A$ যদি
!!a{formula} হয়, তবে $\Box !A$-ও তাই। স্বজ্ঞাতভাবে $\Box !A$ দাবি
করে যে $!A$ \emph{আবশ্যিকভাবে} সত্য, অর্থাৎ যেকোনো সম্ভাব্য জগতে
সত্য। সাধারণত $\Diamond !A$-কে $\lnot \Box \lnot !A$-এর সংক্ষিপ্ত
রূপ ধরা হয় এবং এর অর্থ করা যায় যে $!A$ \emph{সম্ভবত} সত্য। অবশ্য
বিধেয় যুক্তিবিদ্যাতেও মোডালতা যোগ করা যায়।

মোডাল যুক্তিবিদ্যার অর্থতত্ত্ব দিতে ক্রিপকে !!{structure}s ব্যবহার
করা যায়; বস্তুত ক্রিপকে প্রথমে মোডাল যুক্তিবিদ্যার কথাই মনে রেখে এই
অর্থতত্ত্ব নির্মাণ করেছিলেন। আংশিক ক্রমে সীমাবদ্ধ না থেকে আরও সাধারণভাবে
“সম্ভাব্য জগৎ”-এর একটি সেট~$P$ এবং জগৎগুলির মধ্যে একটি দ্বি-স্থানীয়
“অভিগম্যতা” সম্পর্ক $\Atom{R}{x,y}$ নেওয়া হয়। স্বজ্ঞাতভাবে
$\Atom{R}{p,q}$ দাবি করে যে জগৎ~$q$, $p$-এর সঙ্গে সংগত; অর্থাৎ, আমরা
যদি জগৎ~$p$-তে “থাকি”, তবে জগৎটি~$q$-এর মতোও হতে পারত---এই
সম্ভাবনাটি আমাদের বিবেচনায় রাখতে হবে।

ধ্রুপদি যুক্তিবিদ্যার মতো “ব্যাপ্তিগত” যুক্তিবিদ্যার বিপরীতে মোডাল
যুক্তিবিদ্যাকে কখনও কখনও “অভিপ্রায়গত” যুক্তিবিদ্যা বলা হয়। কোনো
ব্যাপ্তিগত যুক্তিবিদ্যার অভিপ্রেত অর্থতত্ত্ব কেবল একটি জগৎ---“প্রকৃত”
জগৎটিরই---উল্লেখ করবে; কিন্তু “অভিপ্রায়গত” যুক্তিবিদ্যার অর্থতত্ত্ব
আরও বিস্তৃত একটি সত্তাতত্ত্বের উপর নির্ভর করে। আবশ্যিকতাকে
!!{structure} দিয়ে গঠন করার পাশাপাশি প্রয়োগ অনুসারে $\Box$ ও
$\Diamond$-এর পুনর্ব্যাখ্যা করে অন্যান্য ভাষাগত নির্মাণকেও
!!{structure} দিয়ে গঠন করা যায়। যেমন:
\begin{enumerate}
\item প্রমাণযোগ্যতা-যুক্তিবিদ্যায় $\Box !A$-কে “$!A$ প্রমাণযোগ্য”
  এবং $\Diamond !A$-কে “$!A$ সঙ্গত” বলে পড়া হয়।
\item জ্ঞানতাত্ত্বিক যুক্তিবিদ্যায় $\Box !A$-কে “আমি $!A$ জানি”
  অথবা “আমি $!A$ বিশ্বাস করি” বলে পড়া যেতে পারে।
\item কালগত যুক্তিবিদ্যায় $\Box !A$-কে “$!A$ সর্বদা সত্য” এবং
  $\Diamond !A$-কে “$!A$ কখনও সত্য” বলে পড়া যায়।
\end{enumerate}

মোডালতা নিয়ে কাজ করার বিধি ও স্বতঃসিদ্ধ দিয়ে যুক্তিবিদ্যাকে সম্প্রসারিত
করতে চাই। যেমন, $\Log{S4}$ ব্যবস্থা সাধারণ প্রস্তাবনামূলক
যুক্তিবিদ্যার স্বতঃসিদ্ধ ও বিধিগুলির সঙ্গে নিচের স্বতঃসিদ্ধগুলি নিয়ে
গঠিত:
\begin{align*}
& \Box (!A \lif !B) \lif (\Box !A \lif \Box !B)\\
& \Box !A \lif !A \\
& \Box !A \lif \Box \Box !A
\intertext{এবং এর সঙ্গে আছে “$!A$ থেকে $\Box !A$ সিদ্ধান্ত করো”
  বিধিটি। $\Log{S5}$ নিচের স্বতঃসিদ্ধটি যোগ করে:}
& \Diamond !A \lif \Box \Diamond !A
\end{align*}
ভিন্ন প্রয়োগের জন্য এই স্বতঃসিদ্ধগুলির ভিন্ন রূপ উপযোগী হতে পারে;
যেমন, সাধারণত S5-কে যৌক্তিক আবশ্যিকতার ধারণাটির চরিত্রায়ক ধরা হয়।
সুবিধার কথা হলো, ক্রিপকে !!{structure}s-এ অভিগম্যতা সম্পর্কটিকে
স্বাভাবিক উপায়ে সীমিত করেই সাধারণত এমন একটি অর্থতত্ত্ব পাওয়া যায়
যার জন্য !!{derivation} ব্যবস্থাটি নির্ভুল ও সম্পূর্ণ। যেমন,
$\Log{S4}$ সেই ক্রিপকে !!{structure}s-গুলির শ্রেণির অনুরূপ যেগুলিতে
অভিগম্যতা সম্পর্কটি প্রতিবিম্ব ও পরিযায়ী। $\Log{S5}$ সেই ক্রিপকে
!!{structure}s-গুলির শ্রেণির অনুরূপ যেগুলিতে অভিগম্যতা সম্পর্কটি
{\em সার্বিক}; অর্থাৎ, প্রত্যেক জগৎ থেকে অন্য প্রত্যেক জগৎ অভিগম্য।
তাই $\Box !A$ সত্য হবে যদি এবং কেবল যদি প্রত্যেক জগতে $!A$ সত্য হয়।

\end{document}
